\documentclass[11pt]{article}
\usepackage[margin=1.2in]{geometry}
\usepackage{amsmath,amssymb,amsthm,booktabs,url}
\usepackage[hidelinks]{hyperref}
\newtheorem{theorem}{Theorem}
\newtheorem{corollary}[theorem]{Corollary}

\title{An Exact Manual Shuffle}
\author{Warren MacEvoy\thanks{Procedure and layout by W.~MacEvoy;
analysis, simulations, and drafting assistance by Claude (Anthropic).}}
\date{August 2026}

\begin{document}
\maketitle

\begin{abstract}
We give a procedure for shuffling a deck of cards by hand, directed by a
die or coins, whose output is \emph{exactly} uniformly distributed over
all $n!$ orderings---not approximately, not asymptotically, but with
zero error in every metric, conditional only on fair independent draws
and one discipline: piles are collected in an order fixed in advance.
The procedure is a single repeating structure that fits on a card table:
deal the deck into a row of bins; collect the bins in predetermined
order; any bin holding more than one card is dealt into the next row
first, by the same rule. It requires no skill, no counting, never shows
a card face, terminates almost surely in a few rows, and a bystander
who observes everything but the card faces---including the random
running time---learns nothing about the result. Run on a computer, the
same rule turns a bare fair-bit generator into an exactly uniform
permutation with output-independent timing. A 52-card deck costs about $121$ card placements using a die and a
coin in a twelve-slot box lid; a purpose-built 36-bin stand with
single-roll resolution of two- and three-card piles reduces this to
about $73$ throws and $55$ placements, with most runs never recursing
at all. The procedure physically realizes the strong uniform time of
Aldous and Diaconis; the proof is a paragraph.
\end{abstract}

\section{The procedure}

Fix a radix $d$ and a randomizer providing independent uniform draws
from $\{1,\dots,d\}$ per card:
\begin{center}
\begin{tabular}{lll}
$d=6$: & one die & (six bins) \\
$d=8$: & three coins read as binary & (eight bins) \\
$d=12$: & one die and one coin & (twelve bins) \\
\end{tabular}
\end{center}

Lay out one \emph{row} of $d$ bin positions, numbered by draw value. The procedure is one rule, applied recursively:

\begin{quote}
\textbf{To shuffle a pile:} deal its cards one at a time, each to the
bin its draw names, all face down. Then collect the bins in a fixed,
predetermined order (say right to left): an empty bin contributes
nothing; a bin holding one card goes onto the output; a bin holding two
or more cards is first shuffled by this same rule, dealt into the next
row below, and its result then goes onto the output.
\end{quote}

Shuffling the deck means applying the rule to the whole deck in
row~one. Each level of recursion is a fresh row of the same $d$ bins,
so the table layout is a small stack of identical rows used top to
bottom; a pair of cards that keeps colliding simply walks one more row
down, and separates there with probability $1 - 1/d$. Seven rows
suffice outside a fraction $\binom{n}{2}/d^{6}$ of runs (about $3\%$
for a 52-card deck at $d=6$, $0.04\%$ at $d=12$); when the rare deeper
collision occurs, reuse a cleared row.

\section{The theorem}

\begin{theorem}
Suppose every draw is independent and uniform on $\{1,\dots,d\}$, and
every collection order is fixed before the corresponding deal. Then the
output of the procedure is exactly uniformly distributed on the $n!$
orderings of the deck. Moreover it is exactly uniform conditional on the
entire observable transcript---every bin size in every row---so an
observer who sees everything except card faces learns nothing about the
result.
\end{theorem}

\begin{proof}
Induction on the pile size $n$; sizes $0$ and $1$ are trivial. Condition
on the size vector $(b_1,\dots,b_d)$ of the top-level deal. Given sizes,
the \emph{sets} of cards in the bins are uniform over all partitions
with those sizes, because the draws are i.i.d.\ uniform; and each bin's
internal order is uniform by the induction hypothesis, independently
across bins. Concatenating in any order fixed independently of the
draws therefore yields a uniformly random ordering of the pile,
conditional on the size vector. Since this holds for every size vector,
it holds unconditionally; and applying the same argument at every level
of the recursion gives uniformity conditional on the full size
transcript. Termination is almost sure: any two cards still sharing a
bin separate at the next row with probability $1 - 1/d$.
\end{proof}

\begin{corollary}
The output permutation is statistically independent of the running
time and of every function of the size transcript: duration, number of
placements, number of re-deals, and rows reached may be disclosed
without any loss of key material. (Uniform conditional on every
transcript value means the joint law factors.) The independence does
not extend to the draw values themselves: an observer of which card
went to which bin reconstructs the permutation entirely, which is why
observation of the draws equals compromise.
\end{corollary}

Equivalently, the procedure sorts the cards by infinite random $d$-ary
strings---by random reals---revealing only as many digits as needed to
break ties; and it is a physical realization of the strong uniform time
of Aldous and Diaconis~\cite{AD86} for riffle-type shuffles: deal by
random digits until all cards' strings are distinct, at which instant
the deck is exactly uniform. Truncating the recursion at a fixed number
of rows, ties broken by deal order, recovers the classical
riffle-shuffle family of Bayer and Diaconis~\cite{BD92} with its small
but nonzero deficits; running the recursion to completion removes the
deficit entirely, which is why no distance tables, entropy accounting,
or pass-count rules appear in this paper. Zero requires no
bookkeeping.

\section{The one discipline}

The hypothesis that collection order is fixed \emph{before} the deal is
load-bearing, and the natural human behavior violates it: pile sizes are
visible, and a tidy operator reaches for the largest pile first. This
biases the result. Already at $n = 3$ cards and $d = 2$, fixed
collection realizes five of the six permutations while
largest-pile-first realizes only three, with total-variation error
$1/2$; the output would always begin with its longest run of
previously-adjacent cards. Independence from the draws is the true
requirement, and predetermination---the same order every row, agreed
before the first roll---is its enforceable form. In multi-party use this
is a security requirement, not hygiene: whoever collects the piles
otherwise holds a channel for biasing the deck while appearing merely
tidy.

\paragraph{Computational form.}
Although stated as a manual scheme, the rule is directly a computer
algorithm, and at $d = 2$ it needs nothing but a fair-bit generator:
stably partition the array by one fresh random bit per element, recurse
on every block of two or more, concatenate. The output is exactly
uniform by the theorem---no modular reduction, no range-rejection loop,
no floating point---and by the corollary its running time is provably
independent of the output, a timing-side-channel guarantee that
implementations of Fisher--Yates must argue for separately. The
expected cost is $n\log_2 n + O(n)$ bits (simulated: $365$ bits at
$n = 52$, $266$ at $n = 40$, against the entropy lower bound
$\log_2 n!$ of $226$ and $159$; overhead factor about $1.6$). The
computational trick itself is folklore---sorting by lazily generated
random reals---but it inherits, unchanged, the exactness proof and the
transcript-independence corollary above.

\section{Multi-party use}

When $N$ parties require a deck none can bias, each rolls its own
randomizer for every card and the bin is the sum modulo $d$ of the
parties' draws (per die; for coins, the XOR of flips). The combined draw
is exactly uniform provided at least one party draws fairly and
independently, so a dishonest coalition of $N-1$ loaded-dice players
gains nothing. Draws for a card are thrown together or pre-committed, so
loaded hardware cannot adapt to honest outcomes; collection needs no
trust, being predetermined and observable; and by the theorem the
procedure may be---should be---conducted with every entitled party
watching everything.

\section{The apparatus}

The radix need not be constant across levels---the theorem's hypotheses
are per-level---and small piles need not recurse at all when their
symmetric group can be sampled directly. Both freedoms condense the
procedure into a purpose-built stand, cuttable on a laser cutter:
\textbf{36 labeled bins} tipped at $60^\circ$ for the first deal
(addressed by two dice, row and column); \textbf{two rows of six
bins} at the right for recursion (one die); and, along the bottom,
\textbf{an action table} for the die faces:
$1,\ 1\mathrm{X},\ 2,\ 2\mathrm{X},\ 3,\ 3\mathrm{X}$.

The operator deals the deck into the 36 bins, then sweeps them in the
stand's manufactured order. A bin of one card goes to the output. Small piles resolve by one roll against the
action table, which labels permutations by \emph{actions} rather than
patterns: the digit names which card (1 = top) moves to the back, and X
says whether to exchange the remaining pair. For a three-card bin this
factors $S_3$ as $3 \times 2 = 6$, one action per die face, exact; for
a two-card bin only the X matters---even faces exchange---exact, three
faces each way. There is no permutation decoding, only two motions:
pull a card to the back, swap or don't. A bin of four resolves in three throws,
deterministically, because $4$ divides $36$: the two dice name a cell,
and each cell of the stand carries a small corner numeral---its index
mod~$4$---so the rig does the arithmetic; the numeral picks the first
card of the pile (exactly uniform, nine cells per residue), and one
further die applies a three-card pattern to the rest. Only a bin of
five or more cards (about one per two 52-card shuffles) is dealt into a
six-bin row and resolved the same way; two rows are ample (cleared rows
are reused), and a majority of 52-card runs never leave the first level
at all. The whole shuffle is thus, in expectation, one deal of the deck
plus roughly a dozen and a half single-roll fix-ups.

Two remarks. First, the stand \emph{manufactures the discipline} of
Section~3: the collection order and the permutation table are cut and
labeled into the artifact, so the theorem's one behavioral hypothesis
is enforced by geometry rather than operator virtue---for a
key-generation apparatus, a property worth having. Second, the
action labels exist for the same reason: pattern strings (``201'' vs.\
``210'') invite systematic misreads, and a systematic misread skews the
six-way distribution; ``move card~2 to the back, exchange'' does not.
What reading error remains is value-independent noise---the operator
sees no faces. A humbler rig
bootstraps the same procedure from four identical slotted shoebox
lids: three tiled side by side form the 36-bin field (two dice: column
and row), and the fourth carries the two recursion rows and the printed
enumeration. A one-page cut template (print four copies at full scale)
accompanies the paper; a single such lid, addressed by a die and a
coin, is by itself the twelve-slot rig of the cost table.

\section{Cost}

Simulated (placements equal draws for the plain recursion; the stand
and shoebox rows count randomizer throws, their single-roll fix-ups
included; ``rows'' is the recursion depth reached):

\begin{center}
\begin{tabular}{lcccc}
\toprule
 & \multicolumn{2}{c}{placements} & \multicolumn{2}{c}{rows} \\
 & mean & 95th pct. & mean & 95th pct. \\
\midrule
$n=52$, $d=6$ (die)          & $157$ & $168$ & $4.8$ & $6$ \\
$n=52$, $d=8$ (coins)        & $138$ & $148$ & $4.3$ & $5$ \\
$n=52$, $d=12$ (die$+$coin)  & $121$ & $129$ & $3.6$ & $5$ \\
$n=52$, shoebox (12 slots, 2/3-fix) & $105$ & $110$ & $2.8$ & $3$ \\
$n=52$, stand (36, $d6$ rows, 4-fix) & $73$  & $79$  & $1.5$ & $2$ \\
$n=40$, $d=12$               & $89$  & $96$  & $3.4$ & $4$ \\
$n=40$, stand                & $53$  & $58$  & $1.2$ & $2$ \\
\bottomrule
\end{tabular}
\end{center}

The recursion re-deals only collisions, never the whole deck, so the
work concentrates tightly around its mean; on the stand, the level-one
deal (52 throws, 52 placements) is in expectation nearly the entire
shuffle---expected placements total $55$, and a majority of runs never
recurse. For comparison, achieving
merely $10^{-3}$ total-variation error with fixed full-deck passes costs
$156$ placements at $n = 52$; the exact procedure with die-and-coin
costs less and errs by nothing.

\section{Cautions}

The randomizers are the entire entropy source; the theorem is
conditional on fair, independent draws---use quality dice or the
multi-party rule. A misdealt card (wrong bin) is harmless
outcome-independent noise, but a collection order improvised after
seeing the piles is not (Section~3). Anyone who observes the draws can
reconstruct the deck: for a secret key, admit exactly the parties
entitled to it, and never reuse a recorded sequence. Simulation code
reproducing the cost table and an empirical uniformity check accompanies
the paper (\texttt{verify\_shuffle.py}, commands \texttt{recursive} and \texttt{stand}).

\begin{thebibliography}{9}
\bibitem{AD86} D.~Aldous and P.~Diaconis,
``Shuffling cards and stopping times,''
\emph{American Mathematical Monthly} \textbf{93} (1986), 333--348.
\bibitem{BD92} D.~Bayer and P.~Diaconis,
``Trailing the dovetail shuffle to its lair,''
\emph{Annals of Applied Probability} \textbf{2} (1992), 294--313.
\bibitem{Schneier99} B.~Schneier,
``The Solitaire encryption algorithm,'' 1999,
\url{https://www.schneier.com/academic/solitaire/}.
\end{thebibliography}

\end{document}
